% META: {"id": "1SPE-VARALEA-EX-029", "chapitre": "1SPE-VARIABLES-ALEATOIRES", "type_objet": "exercice", "capacites": ["1SPE-VARIABLES-ALEATOIRES-C3"], "capacites_codes": ["C3"], "parcours": 1, "competences": ["calculer", "raisonner"], "duree_min": 10, "mode_creation": "ex_nihilo", "sources_inspiration": [], "fichier_tex": "chapitres/1SPE-VARIABLES-ALEATOIRES/exercices/1SPE-VARALEA-EX-029.tex", "corrige_tex": "chapitres/1SPE-VARIABLES-ALEATOIRES/corriges/1SPE-VARALEA-CO-029.tex", "status": "generated"}
% BEGIN-VERIFY
% from sympy import Rational
% # Bernoulli B(0.3): E, V
% p = Rational(3,10)
% E = p; assert E == Rational(3,10)
% V = p*(1-p); assert V == Rational(21,100)
% END-VERIFY
\begin{exercice}{1SPE-VARALEA-EX-029}{1}{10}
Soit $X \sim \mathcal{B}(0{,}3)$ (loi de Bernoulli).

\begin{enumerate}
  \item Dresser le tableau de loi de $X$.
  \item Calculer $E(X)$ et $V(X)$.
\end{enumerate}
\end{exercice}
